\documentclass[tikz,border=12mm]{standalone}
\usepackage[utf8]{inputenc}
\usepackage[T1]{fontenc}
\usepackage[portuguese]{babel}
\usepackage{amsmath,amssymb}
\usepackage{lmodern}
\usepackage{mathpazo}
\usetikzlibrary{arrows.meta,positioning,calc}
\definecolor{ink}{HTML}{263039}
\definecolor{accent}{HTML}{BF5D2E}
\definecolor{sea}{HTML}{275D70}
\definecolor{plum}{HTML}{65577F}
\definecolor{pale}{HTML}{F5F6F4}
\definecolor{line}{HTML}{CED4D5}
\definecolor{cream}{HTML}{F4EDE3}
\begin{document}
\begin{tikzpicture}[
  x=1cm,y=1cm,
  every node/.style={text=ink,font=\fontsize{11.5}{14.5}\selectfont},
  textblock/.style={anchor=north west,align=left,inner sep=0pt,text width=8.45cm},
  heading/.style={anchor=west,font=\sffamily\bfseries\fontsize{15}{18}\selectfont},
  subhead/.style={anchor=west,font=\sffamily\bfseries\fontsize{11.5}{14}\selectfont},
  point/.style={circle,draw=line,fill=white,minimum size=4mm,inner sep=0pt},
  chosen/.style={circle,draw=accent,fill=accent,minimum size=4mm,inner sep=0pt},
  flow/.style={-{Stealth[length=2.4mm]},line width=1.2pt,draw=accent}
]
\useasboundingbox (0,0) rectangle (39.6,27.3);
\node[anchor=west,font=\sffamily\bfseries\fontsize{30}{34}\selectfont]
  at (0,26.65) {FCMP++ da raiz ao pagamento};
\node[anchor=west,font=\fontsize{13}{16}\selectfont]
  at (0,25.55) {Alice conhece a saída que gasta. O nó de Bob conhece a raiz, verifica a prova e reconhece uma imagem já usada.};

\fill[cream,rounded corners=3pt] (0,22.9) rectangle (39.6,24.8);
\node[anchor=west,font=\sffamily\bfseries\fontsize{11}{13}\selectfont,text=accent]
  at (.45,24.25) {RELAÇÃO COMPOSTA};
\node[anchor=west,font=\fontsize{18}{22}\selectfont]
  at (5.6,23.86) {\((O,J,C)\in\mathcal S
    \qquad O=xG+yT
    \qquad I=xJ
    \qquad \widetilde C=C+r_CG\)};
\node[anchor=west,font=\fontsize{10.5}{13}\selectfont]
  at (.45,23.35) {Com Alice ficam o índice do registo, o caminho, as chaves e as máscaras.};
\node[anchor=east,font=\fontsize{10.5}{13}\selectfont]
  at (39.15,23.35) {O nó recebe a raiz \(\rho\), a profundidade \(L\), a mensagem, as imagens e os pseudo-compromissos.};

\foreach \panelx/\panelnum/\paneltitle in {
  0/1/O conjunto e o caminho,
  10.1/2/A prova de pertença,
  20.2/3/A apresentação e SAL,
  30.3/4/A validação do gasto}{
  \fill[pale,rounded corners=3pt] (\panelx,4.75) rectangle (\panelx+9.3,21.95);
  \fill[ink,rounded corners=3pt] (\panelx,20.4) rectangle (\panelx+9.3,21.95);
  \node[heading,text=white] at (\panelx+.4,21.2) {\panelnum\quad \paneltitle};
}
\draw[flow] (9.4,14.0)--(10,14.0);
\draw[flow] (19.5,14.0)--(20.1,14.0);
\draw[flow] (29.6,14.0)--(30.2,14.0);

% O percurso público da construção e a escolha secreta.
\node[textblock] at (.45,19.95)
  {Com 685 saídas, 19 ramos de Selene dão dois ramos de Helios e uma raiz de Selene.};
\node[chosen,label={[font=\small]above:\(\rho\)}] (root) at (4.65,18.0) {};
\node[chosen] (h0) at (3.8,16.65) {};
\node[point] (h1) at (6.9,16.65) {};
\node[point] (s0) at (2.5,15.3) {};
\node[anchor=center,font=\small] at (3.5,15.3) {\(\cdots\)};
\node[chosen] (s1) at (4.65,15.3) {};
\node[point] (s2) at (6.9,15.3) {};
\draw[line,line width=.7pt] (h1)--(root) (s0)--(h0) (s2)--(h1);
\draw[flow] (h0)--(root);
\draw[flow] (s1)--(h0);
\node[anchor=west,text=sea,font=\sffamily\small] at (.5,18) {Selene};
\node[anchor=west,text=plum,font=\sffamily\small] at (.5,16.65) {Helios};
\node[anchor=west,text=sea,font=\sffamily\small] at (.5,15.3) {Selene};
\node[draw=accent,fill=white,rounded corners=2pt,minimum width=8.35cm,minimum height=.9cm]
  (leaf) at (4.65,13.95) {\((O_x,O_y,J_x,J_y,C_x,C_y)\)};
\draw[flow] (leaf.north)--(s1);
\node[textblock,font=\fontsize{10.5}{13}\selectfont] at (.45,13.15)
  {Para a saída 681, os lugares da folha até à raiz são 35, 17 e 0. Só Alice conhece esta escolha.};
\node[subhead] at (.45,11.6) {O nó e a sua apresentação};
\node[textblock] at (.45,11.05) {
  \(\displaystyle Q_j=B_{\mathcal E_j}
     +\langle\boldsymbol v_j,\boldsymbol g_{\mathcal E_j}\rangle\)\\[3mm]
  \(\displaystyle K_j=\langle\boldsymbol v_j,\boldsymbol g_{\mathcal E_j}\rangle
     +\beta_jh_{\mathcal E_j}\)\\[3mm]
  \(\displaystyle B_{\mathcal E_j}+K_j=Q_j+\beta_jh_{\mathcal E_j}\)};
\node[textblock] at (.45,8.5) {
  O vector \(\boldsymbol v_j\) guarda os dados do ramo. O compromisso público \(K_j\) oculta o vector e a máscara \(\beta_j\).};
\node[subhead] at (.45,6.8) {Fecho da raiz por Schnorr};
\node[textblock] at (.45,6.25) {
  \(\displaystyle B_{\mathcal E_L}+K_L-\rho=\beta_Lh_{\mathcal E_L}\)\\[2mm]
  Um único fecho é partilhado pelas entradas.};

% As relações internas da pertença.
\node[textblock] at (10.55,19.95) {
  Os compromissos fixam os ramos, as coordenadas e as funções antes dos desafios.};
\node[draw=sea,fill=white,rounded corners=2pt,minimum width=3.9cm,minimum height=1.3cm,align=center,text=sea]
  at (12.5,18.1) {GBP sobre Selene\\\small camadas ímpares};
\node[draw=plum,fill=white,rounded corners=2pt,minimum width=3.9cm,minimum height=1.3cm,align=center,text=plum]
  at (17.15,18.1) {GBP sobre Helios\\\small camadas pares};
\node[textblock] at (10.55,17.1) {O nó verifica um produto interno por curva.};
\node[subhead] at (10.55,16.1) {A máscara \(W=s_AA\)};
\node[textblock] at (10.55,15.55) {
  \(\displaystyle
    \begin{aligned}
     \operatorname{div}(f_A)
       ={}&\sum_i m_i[A_i]+[-W]\\
          &-(L_A+1)[\mathcal O_\infty].
    \end{aligned}\)\\[3mm]
  A função é fixada antes da secante. Três razões, uma inversa e o teste da curva custam sete produtos secretos.};
\node[subhead] at (10.55,11.9) {Escolher dentro de um ramo};
\node[textblock] at (10.55,11.35) {
  \(\displaystyle \prod_{i=0}^{w-1}(v_i-u)=0\)\\[3mm]
  A lista contém o valor escolhido \(u\). Na folha, seis pesos combinam as coordenadas de cada sextuplo; os dados precedem os pesos.};
\node[subhead] at (10.55,7.8) {Da raiz à saída};
\node[textblock] at (10.55,7.25) {
  A descida desde a raiz fixa os ramos e os sinais dos nós.
  Na folha, obtém-se \((O,J,C)\in\mathcal S\) com as máscaras da apresentação usada por SAL.};

% A interface e a autorização.
\node[textblock] at (20.65,19.95) {
  \(\displaystyle
  \begin{aligned}
   \widetilde O&=O+r_OT,\\
   \widetilde J&=J+r_JU,\\
   \widetilde C&=C+r_CG,\\
   R&=r_JV+r_RT.
  \end{aligned}\)\\[3mm]
  Pertença e SAL recebem esta apresentação comum.};
\node[subhead] at (20.65,15.65) {O produto que liga as respostas};
\node[textblock] at (20.65,15.1) {
  \(\displaystyle w=xr_J\)\\[2mm]
  \(\displaystyle P=xG+r_JV+wU+r_PT\)\\[2mm]
  \(\displaystyle D=P-\widetilde O-R\)};
\node[subhead] at (20.65,12.45) {Aberturas ligadas pelo mesmo desafio};
\node[textblock] at (20.65,11.9) {
  \(\displaystyle
  \begin{aligned}
   \widetilde O&=xG+\bar yT,\\
   D&=wU+\eta T,\\
   I&=x\widetilde J-wU,
  \end{aligned}\)\\[2mm]
  com \(\bar y=y+r_O\) e \(\eta=r_P-\bar y-r_R\).};
\node[subhead] at (20.65,8.8) {A imagem da mesma saída};
\node[textblock] at (20.65,8.25) {
  \(\displaystyle
    I=x(J+r_JU)-(xr_J)U=xJ.
  \)\\[3mm]
  Um SAL por entrada. O desafio vincula a mensagem da transacção, incluindo os pseudo-compromissos.\\[3mm]
  Seis pontos e seis escalares ocupam \textbf{384 bytes}.};

% As propriedades e os testes exteriores.
\node[subhead] at (30.75,19.75) {Pertença e autorização};
\node[textblock] at (30.75,19.2) {
  A raiz admitida fixa o conjunto. A composição extrai o registo \((O,J,C)\) e a abertura \(O=xG+yT\).};
\node[subhead] at (30.75,17.3) {Reconhecer a repetição};
\node[textblock] at (30.75,16.75) {
  A mesma saída produz a mesma imagem \(I=xJ\), pela vinculação das aberturas.\\[2mm]
  As imagens são válidas, não nulas, distintas entre entradas e ainda não usadas.};
\node[subhead] at (30.75,14.2) {Transportar o montante};
\node[textblock] at (30.75,13.65) {
  \(\displaystyle C=zG+aH\)\\[2mm]
  \(\displaystyle\widetilde C=(z+r_C)G+aH\)\\[2mm]
  O coeficiente \(a\) acompanha o mesmo registo até ao pseudo-compromisso.};
\node[subhead] at (30.75,10.8) {Conferir o equilíbrio e os intervalos};
\node[textblock] at (30.75,10.25) {
  \(\displaystyle
    \sum_i\widetilde C_i=\sum_j C_j^{\text{saída}}
                              +\mathrm{taxa}\,H
  \)\\[3mm]
  As provas de intervalo das saídas e os limites das somas asseguram a interpretação inteira dos montantes.};
\node[subhead] at (30.75,7.0) {Ocultação com enunciado fixo};
\node[textblock] at (30.75,6.45) {
  A simulação conserva imagens e pseudo-compromissos, incluindo as correlações impostas pelo equilíbrio.};

\draw[line,line width=.7pt] (0,4.1)--(39.6,4.1);
\node[subhead,text=accent] at (.1,3.55) {MENSAGENS E TAMANHO};
\node[anchor=north west,align=left,text width=19.1cm,inner sep=0pt,font=\fontsize{11.5}{15}\selectfont]
 at (.1,3.0) {
  \(\displaystyle B_{\rm pert}
      =32[\,4(s_{\rm S}+s_{\rm H})+28+
                      2(\log_2m_{\rm S}+\log_2m_{\rm H})\,]\)\\[2mm]
  \(\displaystyle B_{\rm FCMP++}=B_{\rm pert}+480n\).
  \(s\): compromissos por curva; \(m\): dimensão do IPA; \(n\): entradas.\\[2mm]
  Com \(n=1,\ L=3\) são \textbf{4160 bytes} de pertença e \textbf{4640 bytes} no objecto FCMP++.
  Imagens, pseudo-compromissos e provas de intervalo ocupam os seus campos próprios.};
\node[subhead,text=accent] at (20.2,3.55) {DOMÍNIO DAS GARANTIAS};
\node[anchor=north west,align=left,text width=19.1cm,inner sep=0pt,font=\fontsize{10.5}{13}\selectfont]
 at (20.2,3.0) {
  A composição assume extracção das subprovas, dificuldade das relações de logaritmo discreto
  entre as bases e a análise de Fiat--Shamir no modelo de oráculo aleatório.
  Os certificados conservam normalização válida, domínios afins e denominadores admissíveis.\\[2mm]
  A ocultação da realização condicionada exige simulação da mesma distribuição condicionada.
  Os limites dos desafios algébricos aplicam-se a dados previamente fixados.};
\node[anchor=south east,text=sea,font=\sffamily\fontsize{9}{11}\selectfont]
  at (39.6,.05) {ZERO A MONERO \(\,\cdot\,\) FCMP++};
\end{tikzpicture}
\end{document}
